\chapter*{Abstract}
\thispagestyle{previous}

Modern diffusion-based video super-resolution (VSR) can process footage that
runs into minutes, but the field's evaluation instruments cannot: full-reference
metrics require ground truth that deployment videos lack, adjacent-frame
temporal metrics reward spatial smoothness rather than consistency, and
clip-averaged perceptual benchmarks discard the cross-clip drift signal that
long videos accumulate. This thesis characterises these failures empirically and
builds the missing instrument.

First, a systematic failure-mode study on a numerically validated evaluation
environment establishes four structural breakdowns: a content-dependent regime
shift in which two state-of-the-art learned metrics invert method rankings
against human judgement on close-up content; a temporal-scale crossover in which
the conventional tOF metric's winner reverses between adjacent-frame and
long-range lags; a categorical blind spot in which slow colour drift is
invisible to eight widely-used metrics; and a smoother-output bias so strong
that a state-of-the-art consistency benchmark ranks a degraded input above its
own super-resolutions.

Second, the thesis proposes LR-VCC (Long-Range Video Consistency Composite), a
no-reference metric composing seven sub-metrics --- appearance, multi-scale
temporal warping, identity, pairwise and anchored colour stability, semantic
trajectory drift, and colour-trajectory slope --- through a
reliability-weighted softmax-log-mean, in which every sub-metric is gated by an
interpretable regime indicator. Validated on a twelve-family synthetic artefact
benchmark with designed positive and null controls, LR-VCC closes the
colour-drift blind spot, fixes a newly characterised convergence-rewarding
pathology of stability metrics via anchored measurement, and preserves correct
method rankings where existing metrics flip. Applied to real methods (MGLD-VSR,
Upscale-A-Video, FlashVSR), it yields a stable ranking with an actionable
signature: the streaming method wins identity preservation while exhibiting the
worst long-range semantic drift. Hyperparameter-sensitivity and leave-one-out
studies quantify the robustness of each claim.

Third, a causal study localises that drift. A bit-exact position-injection
instrument decomposes rotary position embedding (RoPE) extrapolation into
absolute magnitude and relative geometry, per axis, scored against real ground
truth. Absolute position offsets are quality-free even fifty times beyond the
trained window; relative-distance dilation degrades quality gracefully in time
but steeply in space; real resolution extrapolation is free to twice the
trained extent; and the model's long-video drift is demonstrably not
positional --- exonerating position encoding and redirecting improvement effort
to streaming generation mechanics. The study also demonstrates that
self-consistency does not predict quality, and provides an empirical, per-axis
account of when position interpolation helps and when it hurts.

\textbf{Keywords:} video super-resolution; long video; temporal consistency;
no-reference video quality assessment; rotary position embedding; diffusion
models
\markabstract